\begin{enumerate}
	\item  On a :\quad $\vectt{AB}
	\begin{pmatrix}
		x_\mathrm{B} - x_\mathrm{A}\\
		y_\mathrm{B} - y_\mathrm{A}\\
		z_\mathrm{B} - z_\mathrm{A}\\
	\end{pmatrix}$ \quad soit \quad
	$\vectt{AB}
	\begin{pmatrix}
		2- 1\\
		1-0\\
		-1-3
	\end{pmatrix}$
	\quad et donc, finalement :\quad
	$\vectt{AB}
	\begin{pmatrix}
		1\\
		1\\
		-4\\
	\end{pmatrix}$.

	De même :\quad $\vectt{AC}
	\begin{pmatrix}
		x_\mathrm{C} - x_\mathrm{A}\\
		y_\mathrm{C} - y_\mathrm{A}\\
		z_\mathrm{C} - z_\mathrm{A}\\
	\end{pmatrix}$ \quad soit \quad
	$\vectt{AC}
	\begin{pmatrix}
		1 - 1\\
		1- 0\\
		1-3
	\end{pmatrix}$
	\quad et donc, finalement :\quad
	$\vectt{AC}
	\begin{pmatrix}
		0\\
		1\\
		-2\\
	\end{pmatrix}$.

	Ces coordonnées ne sont pas proportionnelles, donc les vecteurs ne sont pas colinéaires et les points A, B et C ne sont pas alignés.

	Ils définissent donc un plan de l'espace.

	\item Comme le repère est orthonormé, on va utiliser la formule analytique pour calculer les produits scalaires :

	\begin{itemize}[label=\textbullet]
		\item $\vect{n} \cdot \vectt{AB} = 4 \times 1 + 4 \times 1 + 2 \times (-4) = 4 + 4 -8 =0$.\quad $\vect{n}$ et $\vectt{AB}$ sont orthogonaux;
		\item $\vect{n} \cdot \vectt{AC} = 4 \times 0 +4 \times 1 + 2 \times (-2) = 0 + 4 -4 =  0$\quad $\vect{n}$ et $\vectt{AB}$ sont orthogonaux.
	\end{itemize}

	$\vect{n}$ est donc orthogonal à deux vecteurs non colinéaires qui forment une base du plan (ABC), donc il est normal au plan (ABC).



	\item 	Par propriété, on en déduit qu'une équation cartésienne du plan (ABC) est de la forme \\
	$4 \times x + 4 \times y +2  \times z+d=0$ avec $d \in \R$.

	A étant un point du plan (ABC), ses coordonnées doivent vérifier l'équation du plan.

	$\begin{aligned}
		\mathrm{A} \in (\mathrm{ABC})
		&\iff  4 \times 1 + 4\times 0  +2  \times 3 + d = 0\\
		&\iff 4 + 0 + 6 +  d = 0\\
		&\iff d = -10
	\end{aligned}$

	Une équation cartésienne du plan est donc bien \quad$4x + 4y  +2 z -10   = 0$.

	En divisant cette équation par 2 (qui est non nul), on obtient une équation équivalente, donc une équation du plan (ABC) : \quad  $2x + 2y  + z -5   = 0$.

	\item Vérifions que le point H appartient au plan (ABC).

	$2\times x_\mathrm{H}+2 \times y_\mathrm{H}+z_\mathrm{H}-5 = 2 \times 0 + 2\times 2 + 1 -5 = 0+4+1-5=0$

	Donc le point H appartient bien au plan (ABC).

	\item  On a :\quad
	$\vectt{AB}
	\begin{pmatrix}
		1\\
		1\\
		-4\\
	\end{pmatrix}$.

	Et :\quad $\vectt{AH}
	\begin{pmatrix}
		x_\mathrm{H} - x_\mathrm{A}\\
		y_\mathrm{H} - y_\mathrm{A}\\
		z_\mathrm{H} - z_\mathrm{A}\\
	\end{pmatrix}$ \quad soit \quad
	$\vectt{AH}
	\begin{pmatrix}
		0 - 1\\
		2- 0\\
		1-3
	\end{pmatrix}$
	\quad et donc, finalement :\quad
	$\vectt{AH}
	\begin{pmatrix}
		-1\\
		2\\
		-2\\
	\end{pmatrix}$.

	Comme le repère est orthonormé, le produit scalaire peut se calculer avec la formule analytique :
	\quad $\vectt{AB}\cdot\vectt{AH} = 1 \times (-1) +1 \times 2 -4 \times (-2) = -1 +2 +8=9$.


	Comme on est dans un repère orthonormé :

	$\mathrm{AB}=\sqrt{ {x_{\vectt{AB}}}^2  + {y_{\vectt{AB}}}^2  +{z_{\vectt{AB}}}^2 } = \sqrt{ 1^2+ 1^2+(-4)^2}= \sqrt{1+1+16}=\sqrt{18} = 3 \sqrt{2}$.

	et de même :  $\mathrm{AH}=\sqrt{  (-1)^2 +2^2+(-2)^2}= \sqrt{ 1+4+4}= \sqrt{9}=3$.

	L'expression trigonométrique du produit scalaire donne :

	$\begin{aligned}
		\vectt{AB}\cdot\vectt{AH}
		&= \mathrm{AB} \times \mathrm{AH} \times \cos\left(\widehat{\mathrm{BAH}}\right)
	\end{aligned}$

	donc, avec les valeurs connues :\quad $\begin{aligned}[t]
		\vectt{AB}\cdot\vectt{AH}
		&= 3 \sqrt{2} \times 3\times \cos\left(\widehat{\mathrm{BAH}}\right)\\
		&=  9 \sqrt{2} \cos\left(\widehat{\mathrm{BAH}}\right)
	\end{aligned}$

	On a donc deux expressions du même produit scalaire, ce qui donne :

	$ 9= 9 \sqrt{2} \cos\left(\widehat{\mathrm{BAH}}\right) \iff \cos\left(\widehat{\mathrm{BAH}}\right) = \dfrac{1}{\sqrt{2}} = \dfrac{\sqrt{2}}{2} $

	D'où : $\widehat{\mathrm{BAH}} =45 ^\circ$.






	\item La droite $(d)$ orthogonale au plan (ABC)  donc tout vecteur normal au plan (ABC) est un vecteur directeur de la droite $(d)$.

	$(d)$ admet donc le vecteur $\vect{n}$ comme vecteur directeur et passe par le point H.

	Une représentation paramétrique de la droite $(d)$ est donc :

	\[\begin{cases}
		x=0 + 4 \times   t\\
		y=2 + 4 \times t \\
		z=1 +2 \times  t
	\end{cases} \quad \text{avec } t \in \R\]

	\item   Soit S un point de la droite $(d)$.

	Le point H appartient à la fois au plan (ABC) et à  la droite $(d)$  qui est orthogonale au plan (ABC) donc H est le projeté orthogonal du point S sur le plan (ABC).

	La distance entre le point S et  le plan (ABC) est donc égale à la distance SH.

	S est un point de la droite $(d)$ donc il existe un réel $t$ tel que  $\begin{cases}
		x_\mathrm{S} = 4t \\ y_\mathrm{S}=2+4t \\z_\mathrm{S}=1 +2t \end{cases}$.

	On a :

	$\begin{aligned} \mathrm{SH} &= \sqrt{{ \left(  x_\mathrm{H}-x_\mathrm{S} \right)}^2 +  { \left(  y_\mathrm{H}-y_\mathrm{S} \right)}^2+{ \left(  z_\mathrm{H}-z_\mathrm{S} \right)}^2} \\
		& = \sqrt{ (0 - 4t)^2 + (2-2-4t)^2 + (1- 1-2t)^2}\\
		&= \sqrt{16t^2+16t^2+4t^2}\\
		&= \sqrt{36t^2} \\
		&=6 |t| \text{ car une distance est positive}
	\end{aligned}$.


	De plus, la distance entre le point S et le plan (ABC) doit être égale à 6, on a donc :

	$ \begin{aligned} 6 |t| =6  &\iff  |t| = 1\\
		&\iff   t = - 1\text{ ou } t =  1
	\end{aligned}$

	Enfin, l'abscisse $x_{\mathrm{S}}$ doit être positive, on doit donc avoir $4t \geq 0$ soit $t\geq 0$.

	On a donc : $t=1$ et les coordonnés de S sont : $(4;6;3)$.


\end{enumerate}
